\section{数据分析与验证结果}
\label{sec:results}

\subsection{规模与链长}

表\ref{tab:dataset-statistics}给出正式规模。328 条任务链与档案一一对应；总平均链长为 28.37，范围为 5--245。AppWorld 长执行日志将均值显著抬高，因此平均链长不宜作为跨来源质量指标。

\begin{table}[t]
  \centering
  \caption{正式数据规模。5/7 子集只用于说明任务书固定形状的兼容规模。}
  \label{tab:dataset-statistics}
  \small
  \begin{tabular}{lrrrrrr}
    \toprule
    来源 & 源门槛候选 & 发布 & 链长范围 & 5/7 子集 & 档案 & 工具 \\
    \midrule
    $\tau$-bench & 68 & 61 & 5--21 & 25 & 61 & 31 \\
BFCL & 172 & 113 & 5--11 & 46 & 113 & 129 \\
API-Bank & 13 & 7 & 5--6 & 3 & 7 & 14 \\
AppWorld & 147 & 147 & 6--245 & 3 & 147 & 73 \\
合计 & 400 & 328 & 5--245 & 77 & 328 & 247 \\
\bottomrule

  \end{tabular}
\end{table}

$\tau$-bench、BFCL、API-Bank 与 AppWorld 分别发布 61、113、7、147 条。若强制链长只能为 5 或 7，只剩 77 条，即完整数据的 23.5\%；AppWorld 仅保留 3 条。

\subsection{隐私关系与档案来源}

图\ref{fig:data-composition}和表\ref{tab:privacy-statistics}显示档案来源与授权关系。数据包含 1{,}973 个源生字段和 307 个生成字段；允许关系 926 条，禁止关系 25{,}166 条。668 个业务步骤含参考泄露，共新增 1{,}335 个同名敏感字段参数。每条发布任务至少有 4 个敏感字段和 20 个禁止组合。

\begin{figure}[t]
  \centering
  \includegraphics[width=.88\textwidth]{data_composition.pdf}
  \caption{左：四源档案来源分布；右：允许与禁止的字段--工具关系。}
  \label{fig:data-composition}
\end{figure}

\begin{table}[t]
  \centering
  \caption{档案、授权和参考泄露规模。}
  \label{tab:privacy-statistics}
  \scriptsize
  \begin{tabular}{lrrrrrrr}
    \toprule
    来源 & 源生字段 & 生成字段 & 允许 & 禁止 & 泄露步骤 & 触发参数 & 工具参数 \\
    \midrule
    $\tau$-bench & 488 & 0 & 90 & 3,718 & 107 & 214 & 38 \\
BFCL & 545 & 169 & 446 & 16,858 & 253 & 505 & 107 \\
API-Bank & 22 & 27 & 17 & 151 & 20 & 40 & 24 \\
AppWorld & 918 & 111 & 373 & 4,439 & 288 & 576 & 116 \\
合计 & 1,973 & 307 & 926 & 25,166 & 668 & 1,335 & 285 \\
\bottomrule

  \end{tabular}
\end{table}

\subsection{源重放与独立验证}

\begin{table}[t]
  \centering
  \caption{源输入重放和校验结果。AppWorld 输入文件数包含非 Git 任务数据的逐文件哈希。}
  \label{tab:validation-statistics}
  \small
  \begin{tabular}{lrrrr}
    \toprule
    来源 & 重放任务 & 上游文件 & 验证结果 & 错误数 \\
    \midrule
    $\tau$-bench & 61 & 34 & 通过 & 0 \\
BFCL & 113 & 14 & 通过 & 0 \\
API-Bank & 7 & 340 & 通过 & 0 \\
AppWorld & 147 & 751 & 通过 & 0 \\
\bottomrule

  \end{tabular}
\end{table}

表\ref{tab:validation-statistics}表明四个数据包均零错误。328 条功能调用与适配器从锁定输入恢复的参考调用逐项相等。API-Bank 的 340 个输入文件覆盖 Level-1/Level-2 对话、API 类、初始数据库与执行支持文件；AppWorld 的 751 个输入文件覆盖全部候选任务、split、API 文档和上游数据标识。

\begin{table}[t]
  \centering
  \caption{四源环境执行证据。}
  \label{tab:environment-validation}
  \small
  \begin{tabularx}{\textwidth}{Y rrrr Y}
    \toprule
    来源 & 任务数 & 发布执行 & 官方方案 & 最终状态 \\
    \midrule
    $\tau$-bench & 61 & 61 & 0 & 61 & release calls；状态转移对齐 \\
BFCL & 113 & 113 & 0 & 113 & release calls；官方 state checker \\
API-Bank & 7 & 7 & 0 & 0 & release calls；无 final-state oracle \\
AppWorld & 147 & 147 & 147 & 147 & release calls + official solution；原生 task evaluator \\
\midrule
合计 & 328 & 328 & 147 & 321 & 321 release + 147 official \\
\bottomrule

  \end{tabularx}
\end{table}

表\ref{tab:environment-validation}拆分了不同执行对象。四源 328 条发布调用全部在来源特定环境中执行，其中 321 条通过最终状态检查。API-Bank 的 7 条逐步调用原生 API 类，无业务错误，每步通过上游检查器，但因上游未提供整对话最终状态判据而不判定。16 条包含原生业务错误的 $\tau$-bench 参考调用已在发布前过滤。AppWorld 的 147 条发布调用与 147 条官方方案分别在独立环境中执行，均通过任务评测器。

\subsection{语义门禁与人工复核}

正式转换和发布不调用智能体或外部模型。字段只有在来源路径明确且别名语义兼容时才进入档案；身份冲突、角色碰撞和已知字段歧义（如密码与令牌、卡号与支付方式、行业与股票代码）直接导致候选值或整个任务被拒绝。

人工复核从允许边和触发禁止边中按来源确定性抽样，每个来源分别抽取 10 条允许边和 10 条禁止边，共 80 条。复核包隐藏标签和配额，按来源、任务、字段和工具的稳定摘要做全局置换。80 条中 78 条人工判断与确定性标签一致。

\subsection{设计敏感性}

\begin{table}[t]
  \centering
  \caption{完整性、5/7 格式兼容和仅源档案三种派生视图。}
  \label{tab:design-sensitivity}
  \scriptsize
  \begin{tabularx}{\textwidth}{Y rrrrr}
    \toprule
    保留规则 & $\tau$-bench & BFCL & API-Bank & AppWorld & 合计 \\
    \midrule
    完整 reference 主版本 & 61 & 113 & 7 & 147 & 328 \\
仅保留链长5/7 & 25 & 46 & 3 & 3 & 77 \\
仅保留无生成字段 & 61 & 47 & 0 & 39 & 147 \\
\bottomrule

  \end{tabularx}
\end{table}

需要字面满足 5/7 的使用者可从完整数据派生 77 条兼容子集，无需截断参考调用。仅使用无生成字段可得到 147 条仅源档案任务。两种视图都由来源追溯和链长确定性派生，不改变正式文件。
